\section{Założenia Taktyczno-Techniczne}

\paragraph{}
Proces definiowania założeń taktyczno-technicznych stanowi kluczowy etap prac projektowych.
w którym kierownik projektu przeprowadzający wywiad z klientem musi wykazać się nie tylko wszechstronną znajomością branży. 
Lecz również wysokimi kompetencjami mediacyjnymi i krytycznym podejściem do zgłaszanych oczekiwań.

\paragraph{}
Częstym wyzwaniem na tym etapie jest konieczność merytorycznej weryfikacji wyobrażeń zamawiającego.
Ukształtowanych przez materiały promocyjne konkurencji, 
  nierealistyczne oczekiwania połaczone, 
  brakiem znajmości norm oraz ograniczony budżet.

\paragraph{}
Na potrzeby niniejszej pracy przyjęto założenie, 
że zamawiający dysponuje nieograniczonym budżetem projektowym. 
Co zdejmuje z procesu koncepcyjnego standardowe restrykcje kosztowe i potrzeby zatrudnienia do projektu CE \textit{Cost Engineerer}.
Głównym celem projektowym jest stworzenie bezzałogowego systemu powietrznego klasy MALE (\textit{Medium Altitude Long Endurance}) 
dedykowanego do kompleksowej realizacji misji ISR i ISTR. 


W kontekście projektowanej platformy klasy MALE przekłada się to na możliwość bezkompromisowej integracji najnowocześniejszych, 
wielosensorowych modułów ładunku użytecznego.
(zintegrowane głowice EO/IR z podświetlaczami laserowymi, radary SAR o wysokiej 
rozdzielczości oraz systemy SIGINT)[cite: 1, 2], zastosowanie zaawansowanych struktur kompozytowych obniżających sygnaturę radarową oraz 
wdrożenie w pełni redundantnych, zaawansowanych zespołów napędowych i układów awioniki[cite: 1, 2]. 
Taka swoboda koncepcyjna pozwala skoncentrować proces projektowy wyłącznie na maksymalizacji osiągów lotnych –
 w tym uzyskaniu pułapu operacyjnego rzędu $6\,000\text{–}10\,000\text{ m}$, 
 długotrwałości lotu przekraczającej $20\text{–}30\text{ godzin}$ oraz pełnej interoperacyjności zgodnej 
 ze standardami NATO (m.in. STANAG 4671 i STANAG 4586)[cite: 1, 2].



\subsection{Wymagania taktyczno-operacyjne}
Czas lotu
Zasięg
Modulowość
Integraność z sytemami NATO
\subsection{Wymagania techniczne}
Niezbede komponaty do obserwacji takie jak:

\subsection{Przepisy i normy}
STANAG 4702/4703 czy STANAG 4671
\subsection{Podsumowanie}

